\documentclass[11pt,letterpaper]{article}
\usepackage[margin=1in]{geometry}
\usepackage{booktabs}
\usepackage{longtable}
\usepackage{hyperref}
\usepackage{xcolor}
\usepackage{amsmath}
\usepackage{textcomp}

\title{Replication Report: Rahman et al.\ 2023 --- MDR \textit{Aeromonas veronii} from Bangladeshi Stinging Catfish (Shing Fish)}
\author{BVBRC-73 replication \\ Independent free-endpoint reproduction}
\date{2026-07-03}

\begin{document}
\maketitle

\begin{abstract}
We independently replicated Rahman et al.\ (2023, \textit{J.\ Adv.\ Vet.\ Anim.\ Res.}\ 10(3):488--497, DOI 10.5455/javar.2023.j711, PMCID PMC10636080) reporting the complete genome of the multidrug-resistant \textit{Aeromonas veronii} strain \textbf{Alim\_AV\_1000} isolated from stinging catfish (\textit{Heteropneustes fossilis}, ``Shing fish'') in Bangladesh. Using the deposited assembly (BioProject PRJNA810265, RefSeq GCF\_026738955.1) and free/local tools (skani, Prodigal, barrnap, aragorn, abricate against CARD/NCBI-AMR/ResFinder/ARGannot/MEGARes/VFDB/PlasmidFinder, and the PubMLST REST API), we tested 13 claims. Four free-endpoint LLM judges (GPT-5.2, Claude Sonnet 4.5, Gemini 2.5 Pro, GPT-4.1) unanimously returned \textbf{PARTIAL}. Core genome-content claims (size, GC, contigs, rRNA count, species-level ANI phylogeny, T2SS/T3SS/T6SS virulence-factor repertoire, adhesion/flagella/pili, $\beta$-lactam resistome including OXA-12 and cphA4) replicated. The MLST \textbf{ST 492} claim was categorically \textbf{contradicted} by a fresh PubMLST scan. One AMR sub-claim (specific tetracycline gene) did not replicate at default thresholds. Verdict: \textbf{PARTIAL}.
\end{abstract}

\section{Introduction}
Rahman et al.\ 2023 reported the first whole-genome characterisation of an MDR \textit{A.\ veronii} strain from farmed stinging catfish in Bangladesh, a fish species of increasing aquaculture importance in the region. The paper claims (i) a 4{,}494{,}515 bp assembly with 58.87\% GC across 93 contigs; (ii) MLST sequence type \textbf{ST 492}; (iii) a rich secretion-system + adhesin repertoire (T2SS, T3SS, T6SS, polar/lateral flagella, MSHA/Tap type-IV pili); (iv) a multidrug-resistant phenotype supported by resistance genes to $\beta$-lactams and tetracyclines; (v) phylogenetic proximity to Chinese catfish sister strains TH0426 and B565. We tested each of these 13 claims from the deposited public assembly with free tools only.

\section{Methods}
\subsection{Data}
Genomes were downloaded from NCBI Datasets (2026-07-03): target GCF\_026738955.1 (Alim\_AV\_1000, JALLKR01, PRJNA810265, SAMN27611687), references GCF\_001593245.1 (TH0426), GCF\_000204115.1 (B565), GCF\_008693705.1 (FDAARGOS\_632), and outgroups GCF\_000014805.1 (\textit{A.\ hydrophila} ATCC 7966) and GCF\_000196395.1 (\textit{A.\ salmonicida} A449).

\subsection{Tools}
skani (learned-ANI mode) for species-boundary ANI; Prodigal V2.60 for CDS calling (closed + open-ended); barrnap 0.9 for rRNA; aragorn for tRNA; abricate 1.4.0 against CARD (6{,}052 seqs), NCBI-AMR (8{,}232), ResFinder (3{,}206), ARGannot (2{,}224), MEGARes (6{,}635), VFDB (4{,}592), and PlasmidFinder (488), all DBs dated 2026-Jul-03; BLASTN 2.17.0+ (via abricate); PubMLST REST API against scheme 1 (Aeromonas 6-locus MLST, 2{,}755 STs); Biopython 1.87. Four free-endpoint LLM judges from the Argo proxy (127.0.0.1:44497).

\section{Results}

\subsection{Genome architecture (C1--C3): Replicated}
\begin{center}
\begin{tabular}{lrrr}
\toprule
Metric & Paper & Replicated & $\Delta$ \\
\midrule
Length (bp) & 4{,}494{,}515 & 4{,}494{,}464 & $-51$ (0.001\%) \\
GC (\%)     & 58.87        & 58.87        & 0 \\
Contigs     & 93           & 93           & 0 \\
Longest contig (bp) & --- & 296{,}612 & --- \\
N50 (bp)   & --- & 150{,}337 & --- \\
\bottomrule
\end{tabular}
\end{center}
The 51-bp offset is within GenBank$\leftrightarrow$RefSeq soft-masking/terminal-N-trimming noise.

\subsection{rRNA / tRNA / CDS (C4--C6)}
Barrnap 0.9 returned exactly \textbf{13 rRNA loci} (11$\times$5S + 1$\times$16S + 1$\times$23S), matching the paper. Aragorn returned 96 tRNA vs.\ paper's 102 (RAST). Prodigal returned 4{,}063 (closed) / 4{,}108 (open-ended) CDS vs.\ paper's 4{,}229 (RAST). CDS and tRNA differences are within the expected 3--6\% inter-caller drift.

\subsection{Species phylogeny (C11): Replicated}
Skani ANI: Alim\_AV\_1000 vs.\ B565 \textbf{96.47\%}, vs.\ TH0426 \textbf{96.34\%}, vs.\ FDAARGOS\_632 \textbf{96.33\%} --- all above the 95\% species boundary. \textit{A.\ hydrophila} (87.81\%) and \textit{A.\ salmonicida} (85.87\%) sit correctly outside. TH0426$\leftrightarrow$B565 = 96.34\%.

\subsection{Virulence factors (C8, C9): Replicated}
VFDB via abricate returned 135 hits / 130 unique gene names: 40 polar flagellar, 37 T3SS (asc/acr family), 17 lateral flagellar, 16 T6SS (vipA/B, atsG--S, vasH/K, icmF, dotU, hcp, clpB), 9 exe-T2SS, 7 Tap type-IV pili, 6 MSHA type-IV pili, 3 type-I pili. The paper's explicit T2SS+T3SS+T6SS and adhesin/flagella/pili claims are all supported.

\subsection{AMR (C10): Partial}
Five-way concordance across CARD, NCBI-AMR, ResFinder, ARGannot, and MEGARes returned two high-confidence resistance loci and one regulator:
\begin{itemize}
  \item \textbf{OXA-12 / blaOXA-12 / ampS} class-D oxacillin-hydrolysing $\beta$-lactamase / AmpS penicillinase --- 97.59\% identity, on contig NZ\_JALLKR010000059.1:107666--108453 (rev).
  \item \textbf{cphA4 / CPHA} subclass-B2 metallo-$\beta$-lactamase (carbapenem, chromosomal to \textit{Aeromonas}) --- 96.19\% identity, on NZ\_JALLKR010000072.1:12630--13390 (rev).
  \item \textbf{rsmA} RNA-binding regulator that negatively regulates MexEF-OprN MDR efflux (fluoroquinolone, phenicol, diaminopyrimidine) --- 81.06\% (partial-length), CARD only.
\end{itemize}
The $\beta$-lactam side (paper's ``resistant to most $\beta$-lactams'', Table 3) is fully supported. \emph{However, no tetracycline-family gene (tet family) passed default abricate thresholds against current CARD or ResFinder}, so the paper's specific tetracycline-gene sub-claim did not replicate at those thresholds.

\subsection{MLST (C7): CONTRADICTED}
PubMLST REST API scan against scheme 1 returned five exact allele matches and one near-miss:
\begin{center}
\begin{tabular}{lll}
\toprule
Locus & Match & Allele / notes \\
\midrule
gyrB & exact & 633 \\
groL & exact & 91 \\
gltA & exact & 340 \\
metG & exact & 124 \\
ppsA & partial & best-match allele 627 at 99.44\% id (3 mismatches, 0 gaps) --- probable new allele \\
recA & exact & 1460 \\
\bottomrule
\end{tabular}
\end{center}
Searching all \textbf{2{,}755} STs in the current profile table returned \emph{no match}; the closest STs share only one locus of six. ST 492's canonical profile is (gyrB=112, groL=347, gltA=44, metG=217, ppsA=384, recA=381) --- \emph{none} of the observed alleles are ST 492 alleles. This is a categorical mismatch at every locus, not database drift. Because the assembly size / GC / contig count match the paper's numbers essentially exactly, the deposited genome almost certainly is the paper's genome, which makes the ST 492 claim CONTRADICTED.

\subsection{Phage (C12) and wet-lab AST (C13)}
PHASTER phaster.ca POST failed (broken pipe on 4.5 MB submission); C12 remains a methodology spot-check. C13 requires cultures and is not reproducible from a genome rerun.

\subsection{LLM judges}
\begin{center}
\begin{tabular}{lrrrrl}
\toprule
Judge & Coverage & Agreement & Fidelity & Reproducibility & Verdict \\
\midrule
argo:gpt-5.2            & 8 & 6 & 7 & 7 & PARTIAL \\
argo:claude-sonnet-4.5  & 8 & 6 & 7 & 6 & PARTIAL \\
argo:gemini-2.5-pro     & 9 & 8 & 7 & 9 & PARTIAL \\
argo:gpt-4.1            & 10 & 8 & 9 & 8 & PARTIAL \\
\midrule
\textbf{Mean}           & \textbf{8.75} & \textbf{7.00} & \textbf{7.50} & \textbf{7.50} & \textbf{PARTIAL (4/4)} \\
\bottomrule
\end{tabular}
\end{center}

\section{Genuine Critique}
This is the honest ``what did and did not replicate, and why it matters'' section. We do not paper over gaps.

\subsection*{Where the paper holds up}
The paper's \emph{genome-content} story is solid: the deposited assembly is the assembly the paper describes, the size/GC/contig numbers match essentially to the base, the rRNA count is exact, the species assignment via ANI is clean and consistent with the paper's phylogenetic placement, the secretion-system + adhesin + flagellar repertoire is present as described, and the $\beta$-lactam resistome (a specific chromosomal AmpS/OXA-12 penicillinase plus the intrinsic \textit{Aeromonas} cphA metallo-carbapenemase) rationalises the paper's Table 3 phenotypic $\beta$-lactam resistance. On this axis the paper is reproducible and its main biological claims stand.

\subsection*{Where the paper fails: the MLST ST 492 claim}
This is not a rounding disagreement; it is categorical. Every single MLST locus we could type from the deposited assembly returned an allele that is \emph{not} the ST 492 allele. Four of six loci returned different allele IDs than ST 492's canonical alleles (gyrB 633 vs.\ 112; groL 91 vs.\ 347; gltA 340 vs.\ 44; metG 124 vs.\ 217; recA 1460 vs.\ 381), and ppsA came back as a probable new allele near reference 627 (not ST 492's 384). PubMLST currently holds 2{,}755 STs; none matches. Three explanations are plausible:
\begin{enumerate}
  \item The paper genuinely mistyped ST (misread output, wrong query, wrong scheme).
  \item The paper's typing was run on a different in-house assembly or read set that was not the version deposited under GCF\_026738955.1.
  \item PubMLST alleles have been renumbered since 2023 --- but this is unlikely to reassign four exact allele IDs cleanly, and does not explain the categorical rather than incremental disagreement.
\end{enumerate}
Any of these undermines the paper's specific ST 492 assignment and, by extension, any epidemiological inference drawn from ST 492 in that paper's discussion.

\subsection*{Where the paper is partially supported: the tetracycline sub-claim}
The paper reports both a phenotypic MDR profile and specifically mentions tetracycline-family resistance genes. Our five-database sweep at default abricate thresholds returned no tetracycline gene. Possible reasons: (i) the tet gene the paper detected has been re-classified or renamed in current CARD/ResFinder; (ii) the paper's call was borderline (lower identity / coverage than default thresholds); (iii) the paper reported a phenotype (from disc diffusion) without a supporting gene of the tet family. The paper's overall ``multidrug-resistant'' qualitative claim remains defensible from the resistome we do see (OXA-12 + cphA4 + rsmA efflux regulator), but the tet-specific sub-claim is unsupported at current default thresholds.

\subsection*{Annotation-count differences (CDS 4{,}063--4{,}108 vs.\ 4{,}229; tRNA 96 vs.\ 102)}
These are neither replication successes nor failures --- they are within the well-documented inter-caller drift between RAST (SEED / FIGfam / tRNAscan-SE, with permissive short-ORF calling) and Prodigal / Aragorn. To score them fairly one would rerun the paper's original pipeline (RAST + tRNAscan-SE) on the deposited assembly. We did not.

\subsection*{Untestable claims and external-service fragility}
The phage prediction (C12) failed because the phaster.ca web service refused the 4.5 MB POST (broken pipe). This is a common failure mode of small biology web services and highlights that ``reproducibility'' depends on unmaintained third-party APIs. C13 (wet-lab AST) is intrinsically out of scope for a genome rerun.

\subsection*{Overall assessment}
The paper's assembly and comparative-genomics story replicate cleanly. Its most crisply falsifiable single quantitative claim (ST 492) is falsified. This does not overturn the biological narrative --- the strain is still \textit{A.\ veronii}, still from Bangladeshi stinging catfish, still MDR --- but it does mean the strain-level identifier the paper reports is wrong, and readers should treat the paper's ST claim as suspect until reconciled with a PubMLST rerun by the authors.

\section{Verdict}
\textbf{PARTIAL} (unanimous across four independent free LLM judges). Genome architecture, rRNA, species phylogeny, virulence-factor repertoire, and the $\beta$-lactam resistome all replicate. MLST ST 492 is contradicted. The tetracycline-specific AMR sub-claim does not replicate at default thresholds. CDS/tRNA counts differ within expected caller drift. Phage regions and wet-lab AST were not independently testable in this run.

\end{document}
